% ADMD Lecture 2. Compile: latexmk -xelatex lecture02.tex
\documentclass[10pt,aspectratio=169,t]{beamer}
\usetheme{upbminimal}

\course{Application Development for Mobile Devices}
\decklabel{Lecture 2}
\title{Dart and null safety}
\subtitle{The language you will write, and the Flutter toolchain that runs it}
\author{Dinu-Ștefan Rusu}
\institute{FILS, Universitatea Politehnica București}
\date{Autumn 2026}

\begin{document}

\titleframe

\objectivesframe{
  \cell[\faLock]{Null safety, understood}{Read \code{?}, \code{??}, \code{?.}, \code{!} and \code{late}, and know why the compiler can trust them.}
  \cell[\faCode]{Dart you can write today}{Variables, functions, classes, collections, records and pattern matching.}
  \cell[\faIcon{mobile-alt}]{What Flutter draws}{One Dart codebase, one renderer, and the widget catalog at a glance.}
  \cell[\faBolt]{Run the toolchain}{Install, \code{flutter doctor}, hot reload versus hot restart, and a new project's layout.}
}

\divider{Part 1 · Null safety}{The billion-dollar mistake, and its fix}[Why Dart and Kotlin make null part of the type]

\begin{frame}[fragile,eyebrow=Null safety]{The billion-dollar mistake}
  \begin{columns}[T]
    \begin{column}{0.56\textwidth}
\begin{dart}
// Java. Compiles cleanly, no warning.
String name = findUser(42);
System.out.println(name.length());

// findUser is allowed to return null.
// NullPointerException at run time,
// on a user's device, in production.
\end{dart}
    \end{column}
    \begin{column}{0.40\textwidth}
      \lead{Where null came from}{Tony Hoare added the null reference to ALGOL W in 1965 because it was easy to implement. He later called it his "billion-dollar mistake".}
      \muted{Non-nullable by default arrived later, and separately, in Swift, Kotlin, Dart, and TypeScript's strict mode.}
    \end{column}
  \end{columns}
  \begin{takeaway}{Null is not a value problem, it is a type problem.}
    If the type cannot say "maybe absent", the compiler cannot help you.
  \end{takeaway}
  \note{Ask the room what a null pointer exception actually is, in their own words. Most have hit one. Point out this bug is old, well understood, and largely a solved problem in the languages this course uses.}
\end{frame}

\begin{frame}[eyebrow=Null safety]{Type errors are not security bugs}
  \vskip 2mm
  \begin{grid}[3]
    \cell[\faCheckCircle]{Compile-time type errors}{The compiler rejects \code{int x = 'hi'} before it ever ships. Refactoring gets safer, at no runtime cost.}
    \cell[\faExclamationCircle]{Runtime failures}{Null, an out-of-range index, a request that fails. Null safety removes one whole family of these, not all of them.}
    \cell[\faLock]{Security vulnerabilities}{A hard-coded key, a missing check, insecure storage. A fully type-checked app can still be completely insecure.}
  \end{grid}
  \vskip 4mm
  \begin{takeaway}{The one real overlap is memory safety.}
    Dart, Kotlin and Swift have no manual pointer arithmetic, so the classic exploitable bugs of C and C++, buffer overflows and use-after-free, cannot happen there. Security itself stays a design and review discipline; Lecture 10 comes back to it.
  \end{takeaway}
  \note{This slide exists because students conflate a type-safe language with a secure app. Give one concrete example of each column if the room looks unconvinced: a typo caught at compile time, a null from a flaky API, and a hard-coded API key in a public repository.}
\end{frame}

\begin{frame}[fragile,eyebrow=Null safety]{Dart null safety: the operators}
  \begin{columns}[T]
    \begin{column}{0.56\textwidth}
\begin{dart}
// Non-nullable by default:
// String name = null; is an error.
String name = 'Ada';

// Add ? and null becomes allowed.
String? nickname;

// ?. calls only if non-null.
// ?? supplies a fallback.
final shown = nickname?.trim() ?? name;

// ??= assigns only when null.
nickname ??= 'anon';
\end{dart}
    \end{column}
    \begin{column}{0.40\textwidth}
      \lead{Five symbols, and that is most of it}{}
      \begin{itemize}
        \item \code{?} the type may hold null
        \item \code{?.} reads or calls only if not null
        \item \code{??} supplies a fallback value
        \item \code{??=} assigns only when still null
        \item \code{!} asserts non-null, and throws if wrong
      \end{itemize}
    \end{column}
  \end{columns}
  \note{Walk through the code line by line and ask what happens if nickname is null at each point. Emphasize that every one of these operators is a decision the type system forces you to make explicit.}
\end{frame}

\begin{frame}[fragile,eyebrow=Null safety]{What sound null safety means}
  \begin{columns}[T]
    \begin{column}{0.56\textwidth}
\begin{dart}
// Flow analysis promotes a nullable
// to non-nullable within the branch.
String greet(String? name) {
  if (name == null) {
    return 'Hello, stranger';
  }
  return 'Hello, ${name.toUpperCase()}';
  // promoted: no ! needed here
}
\end{dart}
    \end{column}
    \begin{column}{0.40\textwidth}
      \lead{Sound means the guarantee holds}{A non-nullable Dart variable can never hold null at run time. That is a guarantee, not a strong tendency.}
      \vskip 2mm
      \muted{So the compiler stops emitting null checks for it. Sound null safety makes release builds smaller and faster, not only safer.}
    \end{column}
  \end{columns}
  \note{Demo opportunity: delete the null check inside greet and show the analyzer error in the editor, then add ! to make it compile anyway. The lesson is that ! moves the failure from build time to a user's phone.}
\end{frame}

\begin{frame}[fragile,eyebrow=Null safety]{The escape hatches: late and required}
  \begin{columns}[T]
    \begin{column}{0.56\textwidth}
\begin{dart}
// late: promise to assign before use.
class Repo {
  late final Database db;
  Repo() { db = openDatabase(); }
}
// A broken promise throws
// LateInitializationError.

// required: a compile-time obligation.
String card({required String title}) {
  return title;
}
\end{dart}
    \end{column}
    \begin{column}{0.40\textwidth}
      \lead{Two different promises}{\code{late} defers assignment of a non-nullable field. Use it for dependency injection and expensive lazy fields, not to silence the analyzer.}
      \vskip 2mm
      \muted{\code{required} makes a named parameter mandatory. The caller cannot omit it, and the compiler checks that at every call site.}
    \end{column}
  \end{columns}
  \note{Point out that both keywords let you write a non-nullable type where the value is not available yet, one for a field, one for a parameter. Both trade a small promise for a much cleaner type.}
\end{frame}

\begin{frame}[fragile,eyebrow=Null safety]{Null safety: Dart and Kotlin, side by side}
  \begin{columns}[T]
    \begin{column}{0.48\textwidth}
      {\bfseries Dart}
\begin{dart}
String? findEmail(int id) {...}

void notify(int id) {
  final email = findEmail(id);
  if (email == null) return;
  send(email);
  // promoted to String
}
\end{dart}
    \end{column}
    \begin{column}{0.48\textwidth}
      {\bfseries Kotlin}
\begin{kotlin}
fun findEmail(id: Int): String? {}

fun notify(id: Int) {
  val email = findEmail(id)
  if (email == null) return
  send(email)
  // smart cast to String
}
\end{kotlin}
    \end{column}
  \end{columns}
  \vskip 3mm
  \muted{Operators map almost one to one: \code{?.}, \code{??} (Kotlin: \code{?:}), \code{!} (Kotlin: \code{!!}), \code{late final} (Kotlin: \code{lateinit var}, \code{var} only). Learn it once: it transfers straight to Kotlin, and to Swift and Rust after that.}
  \note{This is the slide that pays off for the IoT track students who also write Kotlin. Ask if anyone has hit a NullPointerException in Kotlin interop code; the answer is usually from unannotated Java.}
\end{frame}

\divider{Part 2 · Dart}{The language you will write this semester}[Kotlin alongside it, wherever the comparison helps]

\begin{frame}[fragile,eyebrow=Dart]{Variables: var, final, and const}
  \begin{columns}[T]
    \begin{column}{0.56\textwidth}
\begin{dart}
// var infers its type from the value.
var count = 0;
String title = 'Lecture 2';

// final: assigned once, at run time.
final now = DateTime.now();

// const: computed by the compiler,
// baked into the binary.
const frameBudgetMs = 1000 / 120;

print('$title: $frameBudgetMs ms/frame');
\end{dart}
    \end{column}
    \begin{column}{0.40\textwidth}
      \lead{Read it this way}{}
      \begin{itemize}
        \item \code{var}: still statically typed, only inferred
        \item \code{final}: the binding cannot be reassigned
        \item \code{const}: a compile-time constant
        \item \code{dynamic}: opts out of checking, avoid it
      \end{itemize}
    \end{column}
  \end{columns}
  \note{In Flutter, const widgets are built once and reused across rebuilds, which is a real performance lever, not just a style preference. Save that connection for Lecture 3.}
\end{frame}

\begin{frame}[fragile,eyebrow=Dart]{Functions and named parameters}
  \begin{columns}[T]
    \begin{column}{0.56\textwidth}
\begin{dart}
// Positional; => is a one-expr body.
int add(int a, int b) => a + b;

// Optional positional, with a default.
String tag(String n, [String s = '']) {
  return '$n$s';
}

// Named parameters: the Flutter style.
String card({
  required String title,
  String? subtitle,
}) => title;
\end{dart}
    \end{column}
    \begin{column}{0.40\textwidth}
      \lead{Why named parameters matter}{Every Flutter widget constructor uses them. A call site reads as prose, not as a row of anonymous arguments.}
      \vskip 2mm
      \muted{Functions are first-class values: pass them, store them, return them. That is what a callback like \code{onPressed} really is.}
    \end{column}
  \end{columns}
  \note{Have students read card's call site out loud before showing the definition. Named parameters are the single biggest readability win Flutter borrows from Dart.}
\end{frame}

\begin{frame}[fragile,eyebrow=Dart]{Classes and constructors}
  \begin{columns}[T]
    \begin{column}{0.56\textwidth}
\begin{dart}
class Lecture {
  final int number;
  final String title;
  final String? room;

  // this.x assigns the field directly.
  Lecture(this.number, this.title,
      [this.room]);
}

final l = Lecture(2, 'Languages');
\end{dart}
    \end{column}
    \begin{column}{0.40\textwidth}
      \lead{The constructor is the field list}{\code{this.x} in the parameter list assigns the field with no body needed.}
      \vskip 2mm
      \muted{Mark a field \code{final} and it must be set once, either at declaration or in the constructor, never after.}
    \end{column}
  \end{columns}
  \note{This is the shape every model class in the labs will follow: a handful of final fields, one constructor. Point out room's ? makes it genuinely optional to omit. Current Dart also has primary constructors, which declare the fields in the class header, but the form on this slide is what you will read in every package.}
\end{frame}

\begin{frame}[fragile,eyebrow=Dart]{Getters, equality, and named constructors}
  \begin{columns}[T]
    \begin{column}{0.56\textwidth}
\begin{dart}
class Lecture {
  final int number;
  final String title;

  Lecture(this.number, this.title);

  // Another named constructor.
  Lecture.lab(this.number)
      : title = 'Lab';

  // A getter: computed, no parens.
  String get label => 'L$number';
}
\end{dart}
    \end{column}
    \begin{column}{0.40\textwidth}
      \lead{Equality is identity by default}{Every class gets a default \code{toString}, \code{==}, and \code{hashCode}.}
      \vskip 2mm
      \muted{Two \code{Lecture} objects with the same fields are still not equal. Lecture 5 fixes that with Equatable.}
    \end{column}
  \end{columns}
  \note{Try l1 == l2 on two objects built with the same arguments and watch the room be surprised it prints false. That surprise is exactly the motivation for Equatable later.}
\end{frame}

\begin{frame}[fragile,eyebrow=Dart]{Collections: List, Set, and Map}
  \begin{columns}[T]
    \begin{column}{0.56\textwidth}
\begin{dart}
final topics = ['Dart', 'Kotlin'];
final seen = <String>{};
final points = {'exam': 3, 'lab': 3};

for (final t in topics) {
  seen.add(t);
}
\end{dart}
    \end{column}
    \begin{column}{0.40\textwidth}
      \lead{Three built-ins}{\code{List}, \code{Set} and \code{Map}. The angle brackets, where they appear, name the element types.}
      \vskip 2mm
      \muted{\code{where}, \code{map} and \code{forEach} are lazy \code{Iterable} operations; call \code{.toList()} when you need a \code{List} back.}
    \end{column}
  \end{columns}
  \note{Ask why seen needs the explicit <String>{} while topics does not: type inference reads the list literal's contents, but an empty set literal has nothing to infer from.}
\end{frame}

\begin{frame}[fragile,eyebrow=Dart]{Control flow and collection operators}
  \begin{columns}[T]
    \begin{column}{0.56\textwidth}
\begin{dart}
// Collection-if / collection-for build
// a list inline, with no helper method.
final menu = [
  'Home',
  if (isAdmin) 'Admin',
  for (final t in topics) 'Lab: $t',
];

topics.where((t) => t.length > 4)
    .forEach(print);
\end{dart}
    \end{column}
    \begin{column}{0.40\textwidth}
      \lead{Where you will see this}{Flutter builds a list of children conditionally this way, without a helper function and without a temporary list.}
      \vskip 2mm
      \muted{\code{for}, \code{if} and \code{while} read exactly like Kotlin and Java. Nothing new to learn there.}
    \end{column}
  \end{columns}
  \note{Collection-if and collection-for feel like a small trick the first time you see them, then turn up in nearly every widget build method. Worth a second example if the room looks unsure.}
\end{frame}

\begin{frame}[fragile,eyebrow=Dart]{Records: typed, immutable tuples}
  \begin{columns}[T]
    \begin{column}{0.56\textwidth}
\begin{dart}
// A record: an anonymous tuple type.
(String, int) parseLine(String line) {
  final parts = line.split(',');
  return (parts[0], int.parse(parts[1]));
}

final (name, score) =
    parseLine('Ada,97');
print('$name scored $score');
\end{dart}
    \end{column}
    \begin{column}{0.40\textwidth}
      \lead{One type, many values}{A record bundles values without declaring a class. Two records are equal if their fields are equal.}
      \vskip 2mm
      \muted{Named fields read like a tiny class: \code{(title: String, week: int)}.}
    \end{column}
  \end{columns}
  \note{Records replace the old pattern of returning a small custom class, or worse, a List<Object>, just to hand back two values. The destructuring assignment on the left is new syntax; the next slide names it.}
\end{frame}

\begin{frame}[fragile,eyebrow=Dart]{Pattern matching: switch expressions}
  \begin{columns}[T]
    \begin{column}{0.56\textwidth}
\begin{dart}
String describe(Object v) {
  return switch (v) {
    int n when n < 0 => 'negative',
    int() => 'a whole number',
    String s => 'text: $s',
    _ => 'something else',
  };
}
\end{dart}
    \end{column}
    \begin{column}{0.40\textwidth}
      \lead{A switch that returns a value}{Each arm is a pattern, an optional \code{when} guard, and an expression. \code{\_} is the required wildcard.}
      \vskip 2mm
      \muted{The record from the previous slide can be destructured directly in a \code{case}, matching its shape and binding its fields at once.}
    \end{column}
  \end{columns}
  \note{Compare this to a chain of if-else on runtimeType, which is what students would otherwise reach for. The compiler also warns when a switch expression is not exhaustive.}
\end{frame}

\begin{frame}[fragile,eyebrow=Dart]{Async waits until Lecture 4}
  \lead{One line for now}{Dart handles a value that arrives later, with \code{Future}, \code{async} and \code{await}.}
\begin{dart}
Future<String> fetchName(); // arrives later
\end{dart}
  \begin{takeaway}{Every example in this lecture is synchronous on purpose.}
    The full model, the event loop, error handling and streams, is Lecture 4.
  \end{takeaway}
  \note{Reassure the room that async is coming, not skipped. A Future is just a box that will hold a value later; unpacking it safely is next lecture's entire topic.}
\end{frame}

\begin{frame}[fragile,eyebrow=Dart]{Dart syntax, a quick reference}
  \vskip 2mm
  \begin{tabular}{@{}p{48mm}p{72mm}@{}}
    \toprule
    \thead{Concept} & \thead{Dart} \\
    \midrule
    Variable & \code{var count = 0;} \\
    Immutable binding & \code{final now = DateTime.now();} \\
    Compile-time constant & \code{const budget = 1000 / 120;} \\
    Nullable type & \code{String? name;} \\
    Required named parameter & \code{required this.title} \\
    Typed list literal & \code{<String>['a', 'b']} \\
    \bottomrule
  \end{tabular}
  \vskip 3mm
  \muted{Six lines, and most of the Dart you will read this semester is already here.}
  \note{Use this as a one-slide anchor before moving to Kotlin. Point at each row and ask a student to say out loud what it does, without looking at the previous slides.}
\end{frame}

\begin{frame}[fragile,eyebrow=Dart]{Kotlin and Dart: the same shapes}
  \begin{columns}[T]
    \begin{column}{0.48\textwidth}
      {\bfseries Dart}
\begin{dart}
class User {
  final String name;
  final int age;
  const User(this.name, this.age);
}

final u = User('Ada', 36);
\end{dart}
    \end{column}
    \begin{column}{0.48\textwidth}
      {\bfseries Kotlin}
\begin{kotlin}
data class User(
  val name: String,
  val age: Int,
)

val u = User("Ada", 36)
\end{kotlin}
    \end{column}
  \end{columns}
  \vskip 3mm
  \muted{Kotlin's \code{data class} gives you \code{==}, \code{hashCode}, \code{copy} and \code{toString} for free. In Dart you write them by hand, or generate them: Equatable in Lecture 5, freezed in Lecture 7.}
  \note{The same spectrum from Lecture 1 shows up again here: Kotlin/JVM ships bytecode Android's runtime compiles, Kotlin/Native compiles ahead of time for iOS, the Kotlin Multiplatform path from Lecture 1.}
\end{frame}

\divider{Part 3 · Flutter}{What Flutter actually is}[One Dart codebase, one renderer, and a build mode that changes how the code runs]

\begin{frame}[eyebrow=Flutter]{What Flutter is, and what it draws}
  \begin{columns}[T]
    \begin{column}{0.56\textwidth}
      \begin{itemize}
        \item Google's open-source UI toolkit: one Dart codebase for Android, iOS, web, and desktop.
        \item Flutter does not reuse the platform's own widgets. It ships its own widget library and draws every pixel itself, through the Impeller renderer.
        \item That is why a Flutter app looks identical on both platforms, and why native platform controls are not automatically available.
      \end{itemize}
    \end{column}
    \begin{column}{0.40\textwidth}
      \kv{Your Dart code}{widgets, state, logic}
      \kv{Flutter framework}{layout, animation, gestures}
      \kv{Impeller}{the renderer}
      \kv{GPU}{Metal on iOS, Vulkan or OpenGL on Android}
    \end{column}
  \end{columns}
  \begin{takeaway}{A release build is Dart, compiled ahead of time to native code.}
    There is no interpreter and no JavaScript inside your app.
  \end{takeaway}
  \note{Students who compared Flutter to React Native in Lecture 1 already know the trade-off: Flutter draws its own pixels instead of driving native widgets. This slide is where that fact gets its mechanism.}
\end{frame}

\begin{frame}[eyebrow=Flutter]{The widget catalog, at a glance}
  \vskip 2mm
  \begin{grid}[3]
    \cell[\faLayerGroup]{Layout}{Container, Padding, Row, Column, Stack: how you arrange pixels.}
    \cell[\faKeyboard]{Input}{TextField, Checkbox, ElevatedButton: what the user touches.}
    \cell[\faListUl]{Scrolling}{ListView, GridView: how a long list stays lazy.}
    \cell[\faWindowRestore]{Structure}{Scaffold, AppBar, Drawer: the shell every screen sits in.}
    \cell[\faBell]{Feedback}{SnackBar, AlertDialog, Tooltip: telling the user something happened.}
    \cell[\faToggleOn]{Cupertino}{An iOS-styled set of the same ideas. Lecture 3 goes deeper on all of this.}
  \end{grid}
  \note{This is a map, not a tutorial: point at each cell and name one widget from it that students will meet again in Lecture 3, where layout and composition get a full treatment.}
\end{frame}

\begin{frame}[fragile,eyebrow=Flutter]{Dart compiles two ways: JIT and AOT}
  \begin{columns}[T]
    \begin{column}{0.48\textwidth}
\begin{panel}
{\bfseries Debug: JIT}\par
\code{flutter run}
\vskip 1mm
\begin{itemize}
  \item The Dart VM compiles as it runs
  \item Hot reload swaps in new code in under a second
  \item Assertions on, DevTools attached
\end{itemize}
\end{panel}
    \end{column}
    \begin{column}{0.48\textwidth}
\begin{panel}
{\bfseries Release: AOT}\par
\code{flutter build apk}
\vskip 1mm
\begin{itemize}
  \item Compiled ahead of time to native machine code
  \item No VM ships inside the app
  \item Fast cold start, tree-shaken, no hot reload
\end{itemize}
\end{panel}
    \end{column}
  \end{columns}
  \begin{takeaway}{Never judge performance from a debug build.}
    Measure in profile or release mode: \code{flutter run --profile}.
  \end{takeaway}
  \note{This is the slide that replaces compiled versus interpreted as a language property: same language, two compilers, chosen by build mode. A debug build can be several times slower than release, by design.}
\end{frame}

\begin{frame}[fragile,eyebrow=Flutter]{Hot reload versus hot restart}
  \begin{columns}[T]
    \begin{column}{0.52\textwidth}
      \lead{Hot reload}{\code{r} in the terminal.}
      \begin{enumerate}
        \item You save a \code{.dart} file
        \item Changed source goes to the running VM
        \item The VM re-JITs it in place
        \item Flutter rebuilds the tree; state survives
      \end{enumerate}
    \end{column}
    \begin{column}{0.44\textwidth}
\begin{panel}
{\bfseries What reload cannot do}
\vskip 1mm
\begin{itemize}
  \item Change \code{main()} or start-up code
  \item Change an enum into a class
  \item Fix a compile error first
  \item Run inside a release build at all
\end{itemize}
\end{panel}
    \end{column}
  \end{columns}
  \begin{takeaway}{Hot restart (\code{R}) rebuilds from scratch.}
    All state is lost. Use it whenever \code{initState} or a global variable changed.
  \end{takeaway}
  \note{Ask the room to guess which of the four "cannot do" items is most common in practice. It is usually the compile error one: reload silently does nothing useful until the build succeeds again.}
\end{frame}

\begin{frame}[fragile,eyebrow=Flutter]{Installing the toolchain}
  \begin{columns}[T]
    \begin{column}{0.56\textwidth}
\begin{shell}
# docs.flutter.dev/install

# Android Studio: SDK + emulator
# Xcode (macOS only): iOS builds
# VS Code or IntelliJ, with the
#   Dart and Flutter plugins

flutter doctor
# reports what is missing, per line
\end{shell}
    \end{column}
    \begin{column}{0.40\textwidth}
      \lead{Run it first, always}{\code{flutter doctor} names the exact missing piece. Read that line before asking for help.}
      \vskip 2mm
      \muted{An Android emulator is enough for every lab. A real device is faster, and the only way to test sensors.}
      \vskip 2mm
      \muted{iOS builds need macOS and Xcode; the labs never require them.}
    \end{column}
  \end{columns}
  \note{Lab 2 walks this whole install end to end. Tell students to run flutter doctor before the lab starts, so setup problems surface with time to fix them, not during the two-hour window.}
\end{frame}

\begin{frame}[fragile,eyebrow=Flutter]{Creating and running your first app}
  \begin{columns}[T]
    \begin{column}{0.56\textwidth}
\begin{shell}
flutter create my_app
cd my_app
flutter run
# debug: hot reload r, restart R
flutter run --release
# AOT machine code: what users get
\end{shell}
    \end{column}
    \begin{column}{0.40\textwidth}
      \lead{This week's lab}{Lab 2 walks the install end to end, on Windows, macOS, and Linux.}
      \muted{\code{flutter create} scaffolds a runnable counter app. Read it before you change it.}
    \end{column}
  \end{columns}
  \note{Run flutter create live if time allows. The generated counter app is small enough to read in full, and every student will edit exactly this file in Lab 2.}
\end{frame}

\begin{frame}[eyebrow=Flutter]{The project structure of a new app}
  \begin{tabular}{@{}p{36mm}p{84mm}@{}}
    \toprule
    \thead{Path} & \thead{What it holds} \\
    \midrule
    \code{lib/main.dart} & Entry point: \code{runApp()} and the root widget \\
    \code{pubspec.yaml} & Dependencies, assets, and the SDK constraint \\
    \code{android/} & The native Android project Gradle builds \\
    \code{ios/} & The native Xcode project for iOS builds \\
    \code{test/} & Unit and widget tests, run with \code{flutter test} \\
    \code{build/} & Generated output, never committed to git \\
    \bottomrule
  \end{tabular}
  \vskip 2mm
  \muted{Everything you write lives under \code{lib/}. The rest is generated or native scaffolding you rarely touch.}
  \note{Open a freshly created project side by side with this table. Students consistently ask whether they should edit android/ or ios/ directly; the honest answer for this course is almost never.}
\end{frame}

\begin{frame}[eyebrow=Flutter]{Numbers worth knowing}
  \vskip 4mm
  \begin{grid}[3]
    \bignum{1}{Dart codebase, every platform}
    \bignum{<1 s}{typical hot reload turnaround}
    \bignum{2}{compilers: JIT for debug, AOT for release}
  \end{grid}
  \note{These three numbers are the whole lecture compressed: one codebase, a fast development loop, and a build mode switch that changes how the same source actually runs.}
\end{frame}

\begin{frame}[eyebrow=Recap]{Three things to remember}
  \vskip 2mm
  \begin{enumerate}
    \item Null is part of the type. \muted{\code{?}, \code{??}, \code{?.}, \code{!}, \code{late} and \code{required} remove a whole family of crashes.}
    \item Dart's shapes map onto Kotlin. \muted{Variables, functions, classes, collections, records and patterns: learn the concept once.}
    \item Flutter draws every pixel itself. \muted{JIT while you build for hot reload, AOT for what actually ships.}
  \end{enumerate}
  \vskip 5mm
  \kv{Further reading}{\link{https://dart.dev/language}{dart.dev/language} · \link{https://dart.dev/null-safety}{dart.dev/null-safety} · \link{https://docs.flutter.dev/get-started}{docs.flutter.dev/get-started} · \link{https://kotlinlang.org/docs/null-safety.html}{kotlinlang.org null safety}}
  \note{Close by asking one student to state, in one sentence each, what null safety, records, and Impeller are. If all three land, the lecture did its job.}
\end{frame}

\closingframe{Questions?}{dinu\_stefan.rusu@upb.ro · Microsoft Teams: course channel}

\end{document}
